\clearpage
\section{Diagonalizable Operators}
\subsection{Diagonal Matrices}

\setcounter{thm}{47}
\begin{mydef} [diagonal matrix]
  A \qt{diagonal matrix} is a square matrix that is $0$ everywhere except possibly on \nopagebreak the diagonal.
\end{mydef}

\setcounter{thm}{49}
\begin{mydef} [diagonalizable]
  An operator on $V$ is called \qt{diagonalizable} if the operator has a diagonal matrix with respect to some basis of $V$.
\end{mydef}

\setcounter{thm}{51}
\label{eigenspace}
\begin{mydef} [eigenspace, $E(\lambda, T)$]
  Let $T \in \linmap(V)$ and $\lambda \in \myF$. The \qt{eigenspace} of $T$ corresponding to $\lambda$ is the subspace $E(\lambda, T)$ of $V$ defined by
  \begin{equation}
    E(\lambda, T) :\equiv  \mynull(T-\lambda I) = \{ v\in V \mid Tv = \lambda v\}.
  \end{equation}
\end{mydef}

\setcounter{thm}{53}
\begin{thm} [sum of eigenspaces is a direct sum]
  \label{thm: sum of eigenspaces is a direct sum}
  Suppose $T\in \linmap (V)$ and $\lambda_1, \dots, \lambda_m$ are distinct eigenvalues of $T$. Then
  \begin{equation}
    E(\lambda_1, T) + \cdots + E(\lambda_m, T)
  \end{equation}

  is a direct sum $\left(\text{i.e., } \sum_{i=1}^{m} E(\lambda_i, T) = \bigoplus_{i=1}^{m} E(\lambda_i, T)\right)$.

  Furthermore, if $V$ is finite-dimensional, then
  \begin{equation}
    \begin{aligned}
      \dim E(\lambda_1, T) + \cdots + \dim E(\lambda_m, T)
      & = \dim \left( E(\lambda_1, T)  \oplus \cdots \oplus E(\lambda_m, T) \right) \\
      & \leq \dim V.
    \end{aligned}
  \end{equation}
\end{thm}
\begin{prf}
  Let $\lambda_1, \ddd, \lambda_m$ be distinct eigenvalues. Suppose $v_1 + \cdots + v_m = 0$, where each $v_k \in E(\lambda_k, T)$ for $k = 1, \ddd, m$. This means $\forall k \in \{ 1, \ddd, m\}: v_k = 0$, because eigenvectors corresponding to distinct eigenvalues are linearly independent by \ref{thm: linearly independent eigenvectors}. Thus, $E(\lambda_1, T) + \cdots + E(\lambda_m, T)$ is a direct sum by \ref{thm: condition for a direct sum}.
  By \ref{thm: a sum is a direct sum if and only if the dimensions add up}, the dimensions of a direct sum add up, and by \ref{thm: dimension of a subspace}, the dimension of a subspace is always $\leq$ to the space itself (if it is \fd). Therefore,
  \begin{equation}
    \dim \left( \underbrace{E(\lambda_1, T)  \oplus \cdots \oplus E(\lambda_m, T)}_{\text{subspace of $V$}} \right)
    =  \dim E(\lambda_1, T) + \cdots + \dim E(\lambda_m, T) \leq \dim V.
  \end{equation}
  \vspace{-1em}
\end{prf}

\subsection{Conditions for Diagonalizability}
\setcounter{thm}{54}
\begin{thm} [conditions equivalent to diagonalizability]
  \label{thm: conditions equivalent to diagonalizability}
  Let $\lambda_1, \dots,\lambda_m$ denote the distinct (i.e., $\lambda_j \neq \lambda_k$) eigenvalues of $T\in \linmap (V)$. Then the following are equivalent:
  \begin{enumerate}[label=\textbf{(\alph*)}]
    \item $T$ is diagonalizable.
    \item $V$ has a basis consisting of eigenvectors of $T$.
    \item $V=E(\lambda_1, T) \oplus \cdots \oplus E(\lambda_m, T).$
    \item $\dim V = \dim E(\lambda_1, T) + \cdots + \dim E(\lambda_m, T)$.
  \end{enumerate}
\end{thm}
\begin{prf} We prove (a) $\iff$ (b), (b) $\implies$ (c), (c) $\implies$ (d), and then (d) $\implies$ (b). Let $\lambda_1, \ddd, \lambda_m$ be distinct eigenvalues of $T \in \linmap (V)$ for the whole proof.

  \StepOne Suppose (a) holds. The matrix of $T$ looks as follows:
  \begin{equation}
    \mmatrix(T) =
    \left( {\begin{array}{ccc}
        \lambda_1 &         & 0 \\
        &  \ddots &    \\
        0      &         & \lambda_m
    \end{array} } \right)
  \end{equation}

  with respect to a basis $v_1, \ddd, v_m$. This is the case $\iff Tv_k = \lambda_k v_k$ for $k = 1, \ddd, m$. Thus, (a) and (b) are equivalent, i.e., (a) $\iff$ (b).

  \StepTwo Now, suppose (b) holds; thus, $V$ has a basis consisting of eigenvectors of $T$. Hence, every vector in $V$ is a linear combination of such basis, which implies
  \begin{equation}
    V = E(\lambda_1, T) + \cdots + E(\lambda_m, T).
  \end{equation}

  By \ref{thm: sum of eigenspaces is a direct sum}, the distinctness of all $\lambda_j$'s implies that this is a direct sum $\left(\text{i.e., } V=\bigoplus_{j=1}^{m}  E(\lambda_j, T) \right)$. This proves that (b) $\implies$ (c).

  \StepThree (c) $\implies$ (d) follows immediately from \ref{thm: a sum is a direct sum if and only if the dimensions add up}, which states that the dimensions of a direct sum add up $\left(\text{i.e., } \dim \left(\bigoplus_{j=1}^m E(\lambda_j, T)\right) = \sum_{j=1}^{m} \dim V_j \right)$.

  \StepFour Finally, suppose (d) holds; thus,
  \begin{equation}
    \dim V = \dim E(\lambda_1, T) + \cdots + \dim E(\lambda_m, T).
  \end{equation}

  Choose a basis of each $E(\lambda_m, T)$; put all these together to form a list
  \begin{equation}
    \label{eq: list of vectors v}
    v_1, \ddd, v_n
  \end{equation}
  of eigenvectors of $T$, where $n = \dim V$. We want to show that \eqref{eq: list of vectors v} forms a basis of $V$. To this end, suppose for some $a_1, \ddd, a_n \in \myF$:
  \begin{equation}
    \label{eq: sum of vs}
    a_1 v_1 + \cdots + a_n v_n = 0.
  \end{equation}

  Now, for every $k = 1, \ddd, m$, let
  \begin{equation}
    u_k := \hspace{-1.5em} \sum_{\substack{v_j \in E(\lambda_k, T), \text{ such that} \\ \text{$v_j$ is a basis vector of $E(\lambda_k, T)$}}} \hspace{-1.5em} (a_j v_j)
  \end{equation}

  denote the $k$-th sum of the terms $a_jv_j$, where the $v_j$'s belong to the basis of $E(\lambda_k, T)$. This means that $u_k \in E(\lambda_k, T)$ is an eigenvector with eigenvalue $\lambda_k$.
  Therefore, for every $k = 1, \ddd, m$, we have
  \begin{equation}
    u_k \in E(\lambda_k, T)
  \end{equation}

  and
  \begin{equation}
    u_1 + \cdots + u_m = 0,
  \end{equation}

  as this is the same sum we had earlier \eqref{eq: sum of vs}. By \ref{thm: linearly independent eigenvectors},
  every list of eigenvectors of $T$ corresponding to distinct eigenvalues of $T$ is linearly independent. This forces
  \begin{equation}
    u_k = 0 \quad \forall k \in \{1, \ddd, m \}.
  \end{equation}

  This implies that all the $a$'s must be $0$, because the $v$'s are bases of the $E$'s (eigenspaces), and their corresponding sums (the $u$'s) all equal $0$. Therefore, $v_1, \ldots, v_n$ is a linearly independent list, and hence, a basis of $V$, by \ref{thm: linearly independent list of the right length is a basis} ($n = \dim V$). Now, we have shown that (d) $\implies$ (b), completing the proof.
\end{prf}

\setcounter{thm}{57}
\begin{thm} [enough eigenvalues implies diagonalizability]
  \label{thm: enough eigenvalues implies diagonalizability}
  If $T\in \linmap(V)$ has $\dim V$ distinct eigenvalues, then $T$ is diagonalizable.
\end{thm}
\begin{prf}
  Let $\lambda_1, \ldots, \lambda_{\dim V}$ be the eigenvalues of the eigenvectors $v_1, \ldots, v_{\dim V}$, which are linearly independent by \autoref{thm: linearly independent eigenvectors}. So, we have a basis consisting of $\dim V$ eigenvectors. Hence, $T$ is diagonalizable by \ref{thm: conditions equivalent to diagonalizability}.
\end{prf}

\setcounter{thm}{61}
% 5.62
\begin{thm}[necessary and sufficient condition for diagonalizability]
  \label{thm: necessary and sufficient condition for diagonalizability}
  Let $T\in \linmap (V)$, $\dim V \neq \infty$. Then $T$ is diagonalizable $\iff$ the minimal polynomial of $T$ equals
  \begin{equation}
    (z-\lambda_1) \cdots (z-\lambda_m),
  \end{equation}

  where $\lambda_1, \dots, \lambda_m \in \myF \myand \lambda_i \neq  \lambda_j $ for $ i,j \in \{1, \ddd, m \}$.
\end{thm}
\begin{prf} For this proof, let $T\in \linmap (V)$ be finite-dimensional ($\dim V \neq \infty$).

  \Rightarrowdirection Suppose $T$ is diagonalizable. Thus, there exists a basis $v_1, \ddd, v_n$ of $V$ consisting of eigenvectors of $T$. Let $\lambda_1, \ddd, \lambda_m$ be the distinct eigenvalues of $T$. Therefore, each $v_j$ is an eigenvector corresponding to some $\lambda_k$, so $(T-\lambda_k I)v_j = 0.$
  Thus,
  \begin{equation}
    (T-\lambda_1 I) \cdots (T-\lambda_m I)v_j =0,
  \end{equation}

  which implies the minimal polynomial equals $(z-\lambda_1)\cdots(z-\lambda_m).$ (Why?)

  \Leftarrowdirection Suppose the minimal polynomial of $T$ is $(z-\lambda_1) \cdots (z-\lambda_m)$ for distinct scalars $\lambda_1, \ddd, \lambda_m \in \myF$.
  Write
  \begin{equation}
    \label{eq: (t-l1 I)***(t-l_{m}I)=0}
    (T-\lambda_1 I) \cdots (T-\lambda_{m-1} I) (T-\lambda_m I)=0.
  \end{equation}

%  Hence, if we drop one factor $((T-\lambda_1 I) \cdots (T-\lambda_{m-1} I) \xcancel{(T-\lambda_m I)})$ we get
%  \begin{equation}
%
%    (T-\lambda_1 I) \cdots (T-\lambda_{m-1} I)=0.
%  \end{equation}

  Suppose $m=1$. Then $T-\lambda_1 I = 0 \iff T=\lambda_1 I$, so $T$ is diagonalizable. This is our base case for induction over $m \in \nat$.
  Now, suppose $m>1$ and that the desired result holds for all values smaller than $m$. Like in the proof of Theorem \ref{thm: necessary and sufficient condition to have an upper-triangular-matrix}, we observe that $\myrange (T-\lambda_m I)$ is invariant under any polynomial of $T$ (using \ref{thm: null space and range of p(T) are invariant under T} with $p(z) := z-\lambda_m$). Now, define
  \begin{equation}
    \label{eq: 2nd definition of U}
    U := \myrange(T - \lambda_m I).
  \end{equation}

  This means that $\left . T \right |_{U} $ is an operator on $U$.
  If $u \in \myrange (T-\lambda_m I)$, then $u = (T-\lambda_m I)v$ for some $v \in V$. Together with \eqref{eq: (t-l1 I)***(t-l_{m}I)=0}, we have that
  \begin{align}
    \label{eq: equation for u and v}
    (T-\lambda_1 I) \cdots (T- \lambda_{m-1} I)u
    &=(T-\lambda_1 I) \cdots (T-\lambda_m I)v \\
    &=0. % \mytext{(since $v$ consists of $v_j$'s)}
  \end{align}

  Since $u \in U$ was arbitrary, $q(z) := (z-\lambda_1)\cdots (z-\lambda_{m-1})$ is a polynomial \st $q(T)(u)=0$ for every $u \in U$. By \ref{thm: every zero polynomial is a multiple of the minimal polynomial}, $q(z)$ is a polynomial multiple of the minimal polynomial of $\left . T \right |_{U}$.
  By our induction hypothesis, there is a basis of $U$ consisting of eigenvectors of $\left . T \right |_U$ and therefore, as well as eigenvectors of $T$.
  Suppose
  \begin{equation}
    x\in U \cap \mynull (T-\lambda_m I) = U \cap E(\lambda_m, T).
  \end{equation}

  Hence, $Tx = \lambda_m x$, because $Tx -\lambda_m Ix=(T-\lambda_m I)x = 0$. Since $x \in U$, equation \eqref{eq: equation for u and v} implies that
  \begin{equation}
    \begin{aligned}
      0 & = (T-\lambda_1 I) \cdots (T- \lambda_{m-1} I)x \\
      & = (\lambda_{m}I-\lambda_1 I) \cdots (\lambda_{m}I- \lambda_{m-1} I)x \\
      & = (\lambda_{m}-\lambda_1 ) \cdots (\lambda_{m}- \lambda_{m-1} )x. \\
    \end{aligned}
  \end{equation}

  Because $\lambda_i \neq \lambda_j$ for all $i \neq j$ (in particular, $\lambda_m - \lambda_k \neq 0$ for all $k < m$), this can only happen if $x=0$. Therefore,
  \begin{equation}
    U \cap \mynull (T-\lambda_m I) = \{0\}.
  \end{equation}

  Thus, $U + \mynull (T-\lambda_m I)$ is a direct sum by \ref{thm: sum and intersection of two subspaces}. That is,
  \begin{equation}
    U + \mynull (T-\lambda_m I) = U \oplus \mynull (T-\lambda_m I).
  \end{equation}

  Since $(T-\lambda_m I) \in \linmap (V)$, we can use the Rank-Nullity Theorem \ref{rank-nullity-theorem} to conclude that
  \begin{equation}
    \dim V = \dim U + \dim \mynull (T-\lambda_m I).
  \end{equation}

  \ref{thm: a sum is a direct sum if and only if the dimensions add up} states that a sum is a direct sum $\iff$ the dimensions add up to the dimension of the sum. This means that we can rewrite the equation above to:
  \begin{equation}
    \dim V = \dim (U \oplus \mynull (T-\lambda_m I)).
  \end{equation}

  According to \ref{thm: subspace of full dimension equals the whole space}, together with the fact that $U \oplus \mynull (T-\lambda_m I) \subseteq V$, this implies
  \begin{equation}
    V=U \oplus \mynull (T-\lambda_m I) = U \oplus E(\lambda_m, T).
  \end{equation}

  Every vector in $\mynull (T-\lambda_m I)$ is an eigenvector of $T$ with eigenvalue $\lambda_m$. Earlier in this proof, we saw that there is a basis of $U=\myrange(T-\lambda_mI)$ consisting of eigenvectors of $T$ (using our induction hypothesis on this smaller subspace of $V$). Adjoining to that basis a basis of $E(\lambda_m, T) = \mynull(T-\lambda_mI)$ gives a basis of $V$ consisting of eigenvectors of $T$. The matrix $\mmatrix(T)$ of $T$ with respect to this basis is a diagonal matrix, as desired, using \ref{thm: conditions equivalent to diagonalizability} (according to that theorem, $V$ has a basis consisting of eigenvectors of $T \in \linmap(V) \iff$ it is diagonalizable). As a side note, this theorem states that $V = E(\lambda_1, T) \oplus \cdots \oplus E(\lambda_m, T)$, which correlates with $V=U\oplus \mynull (T-\lambda_m I)$.
\end{prf}

\setcounter{thm}{64}
%5.65
\begin{thm}[restriction of diagonalizable operator to invariant subspace]
  \label{thm: restriction of diagonalizable operator to invariant subspace}
  Let $T\in \linmap(V)$. Suppose $T$ is diagonalizable and $U$ is a subspace of $V$ that is invariant under $T$. Then
  \begin{equation}
    \left.T\right|_U \text{ is a diagonalizable operator on } U.
  \end{equation}
\end{thm}
\begin{prf}
  %  Diagonazability of $T$ $\mathsmaller{\overset{\text{\ref{thm: necessary and sufficient condition for diagonalizability}}}{\iff}}$ the minimal polynomial of $T$ equals
  Diagonalizability of $T \iff_{(\ref{thm: necessary and sufficient condition for diagonalizability})}$ the minimal polynomial of $T$ equals
  \begin{equation}
    (z-\lambda_1)\cdots(z-\lambda_m) \mytext{for} \lambda_j \neq  \lambda_k.
  \end{equation}

  By \ref{thm: minimal polynomial of a restriction operator}, the minimal polynomial of $T$ is a polynomial multiple of the minimal polynomial of $\left.T\right|_U$.
  Hence, the minimal polynomial of $\left.T\right|_U$ has the form required by \ref{thm: necessary and sufficient condition for diagonalizability}, which shows that $\left.T\right|_U$ is diagonalizable. It consists of any factor of the form $(z-\lambda_k)$.
\end{prf}

% 5.66
\setcounter{thm}{65}
\begin{mydef}[Gershgorin disk]
  \label{def: gershgorin disk}
  Let $A \in \myF^{n,n}$. For $j \in \{1, \dots, n\}$, let $r_j := \sum_{k \neq j} |A_{j,k}|$ denote the sum of the absolute values of the off-diagonal entries in the $j$-th row of $A$. The $j$-th \qt{Gershgorin disk} is the closed disk in the complex plane defined by
  \begin{equation}
    D(A_{j,j}, r_j) := \{ z \in \compl \mid |z - A_{j,j}| \leq r_j \}.
  \end{equation}
\end{mydef}

% 5.67
\begin{thm}[Gershgorin disk theorem]
  \label{thm: gershgorin disk theorem}
  Let $A \in \myF^{n,n}$. Every eigenvalue of $A$ is contained in some Gershgorin disk of $A$. In other words, if $\lambda \in \myF$ is an eigenvalue of $A$, then
  \begin{equation}
    \lambda \in \bigcup_{j=1}^n D(A_{j,j}, r_j).
  \end{equation}
\end{thm}
\begin{prf}
  Let $x = (x_1, \dots, x_n)^\top \neq 0$ be an eigenvector of $A$ corresponding to $\lambda$, so $Ax = \lambda x$.
  Choose an index $j \in \{1, \dots, n\}$ such that $|x_j| = \max \{|x_1|, \dots, |x_n|\} > 0$.
  The $j$-th row of the equation $Ax = \lambda x$ gives $\sum_{k=1}^n A_{j,k} x_k = \lambda x_j$.
  Rearranging gives
  \begin{equation}
    (\lambda - A_{j,j}) x_j = \sum_{k \neq j} A_{j,k} x_k.
  \end{equation}
  Taking absolute values and using the triangle inequality:
  \begin{equation}
    |\lambda - A_{j,j}| \, |x_j| \leq \sum_{k \neq j} |A_{j,k}| \, |x_k| \leq \sum_{k \neq j} |A_{j,k}| \, |x_j| = r_j |x_j|.
  \end{equation}
  Dividing by $|x_j| > 0$ yields $|\lambda - A_{j,j}| \leq r_j$, so $\lambda \in D(A_{j,j}, r_j)$.
\end{prf}